% titel: Blatt 10, A5: Tabelle pH/pOH/c(H+)/c(OH-) + Kw(0°C)
% Teil II: der durchgerechnete Weg.

\subsection*{a) Tabelle auffüllen (25\,°C)}
Gegeben c(H$^+$)=7.5$\times10^{-3}$:
\[
  c(\mathrm{OH^-}) = \frac{10^{-14}}{7.5\cdot10^{-3}} \approx 1.33\cdot10^{-12},
  \quad pH = 2.12,\quad pOH = 11.88,\quad \text{sauer}.
\]
Gegeben pH=5.70:
\[
  c(\mathrm{H^+}) = 10^{-5.70} \approx 2.0\cdot10^{-6},
  \quad c(\mathrm{OH^-}) = \frac{10^{-14}}{2.0\cdot10^{-6}} \approx 5.0\cdot10^{-9},
  \quad pOH = 8.30,\quad \text{sauer}.
\]

\subsection*{b) Neutrale Lösung bei 0\,°C}
$K_w = 1.2\cdot10^{-15}$.
\[
  c(\mathrm{H^+}) = c(\mathrm{OH^-}) = \sqrt{K_w}
   = \sqrt{1.2\cdot10^{-15}} \approx 3.5\cdot10^{-8}\ mol\,L^{-1}.
\]

\subsection*{c) Faktor bei pH-Änderung}
$c(\mathrm{H^+})=10^{-pH}$, also Änderung um $\Delta pH$ gibt Faktor
$10^{-\Delta pH}$:
\begin{itemize}
  \item $\Delta pH = 2.00 \Rightarrow c(\mathrm{H^+})$ um den Faktor
        $10^{-2}$ (also $1/100$, bzw. um $10^2$ ändert sich der Stofficid).
  \item $\Delta pH = 0.50 \Rightarrow$ Faktor $10^{-0.5}\approx 3.2$
        (bzw. $1/3.2$).
\end{itemize}

\begin{ergebnis}
a) Zeile 1: OH$^-\approx1.3e{-}12$, pH$=$2.12, pOH$=$11.88, sauer. Zeile 2:
H$^+\approx2.0e{-}6$, OH$^-\approx5.0e{-}9$, pOH$=$8.30, sauer.
b) $c(\mathrm{H^+})=c(\mathrm{OH^-})\approx3.5\times10^{-8}$ mol/L (0\,°C).
c) Faktor $10^{-2}$ ($\Delta pH{=}2$) bzw.\ $10^{-0.5}\approx0.32$
($\Delta pH{=}0.5$).
\end{ergebnis}

\begin{kontrolle}
In beiden Zeilen $pH+pOH=14$ erfüllt, Produkt $=10^{-14}$. Bei 0\,°C ist
$K_w$ kleiner $\Rightarrow$ neutrales $c$ nur $\approx3.5e{-}8$ (kleiner als
bei 25\,°C), plausibel. 
\end{kontrolle}

\begin{merksatz}
Vier Größen hängen über $K_w=c(\mathrm{H^+})c(\mathrm{OH^-})$ und
$pH+pOH=14$ zusammen. $K_w$ temperaturabhängig; pH-Faktor ist eine
Zehnerpotenz.
\end{merksatz}
